% Word count aim: 1500-2000
\chapter{Introduction}
\label{cha:intro}

% Put in an illustrative picture
\section{Motivation}
Many aspects of our day-to-day lives revolves around finding
optimality in problems that we can have no comprehensible insight
about. It can be observed in industrial processes of ! TODO!. These
types of problems are said to be \acrfull{bbo} problems.

We do not get much tools to interact with them save from a
cost/objective function $\mathbf{f}: \mathcal{X} \mapsto \mathcal{S}$
(the `oracle') that provides an abstract representation of the
problem, where $\mathcal{X}$ is the search space, $\mathcal{S}$ the
evaluation space !TODO: find concrete refs!. Typically, $\mathcal{S}
\subseteq R$. The task is to minimize or maximize the solutions
produced by this objective function, i.e., find a solution that Much
study is done in optimizing this cost function without the liberty of
classic mathematical tools such as continuity, smoothness and
convexity (these methods warrant the oracle to return additional
  information, such as gradients, apart from the evaluation value of an
input \cite{liu_primer_2020}).

\subsection{Conventional Solutions}
% A brief history of how this problem came to be
Classical analysis of solving \acrshort{bbo} problems are concerned
with the so-called \acrfull{bbc}, that is, deriving the runtime
bounds for an algorithm to sample a (global) optimum, depending only
on the problem size $\mathbf{n}$. The reasoning behind this is
palpable; the \textit{computational} complexity of evaluating the
objective function for such \acrshort{bbo} systems are usually expensive.

A trivial example is Hyperparameter Tuning in \acrfull{ml} models
!CITE!. The optimal hyperparameters could be thought of as the
\textit{optima} for the \acrfull{bbo} problem of the tuning that maximizes
their validation loss\footnote{`how well the model fits the
underlying probability distribution of the data inputs and outputs'}.

\newpage
% Give context to the reader about the objectives
\section{Aims and Objectives}
% \subsection{General Aims}
% % TODO
% \subsection{Theoretical Exploration}
% % TODO
% \subsection{Supplementary Software Development} % State hypotheses here
% % TODO
% \subsection{Experimental Analysis} % State hypotheses here
% % TODO
\begin{longtable}{||p{0.5\textwidth}|p{0.5\textwidth}||}
\hline
\textbf{Broader Aims} & \textbf{Specific Aims} \\
\hline
\endhead
\hline
\textbf{Aim 1. } Derive a viable relationship between the canonical measure of \acrshort{bbo} algorithms and regret.
&  
\begin{enumerate}
    \item \textit{Aim 1.1. } Seek out the standard measures used in \acrshort{bbo} theory. 
    \item \textit{Aim 1.2. }
    \item \textit{Aim 1.2. }
    \item \textit{Aim 1.2. }
    \item \textit{Aim 1.2. }
\end{enumerate}
\\
\hline
\textbf{Aim 2. } Derive a viable relationship between the canonical measure of \acrshort{bbo} algorithms and regret.
&  
\begin{enumerate}
    \item \textit{Aim 2.1. } Seek out the standard measures used in \acrshort{bbo} theory. 
    \item \textit{Aim 2.2. }
    \item \textit{Aim 2.2. }
    \item \textit{Aim 2.2. }
    \item \textit{Aim 2.2. }
\end{enumerate}
\\
\hline
\textbf{Aim 3. } Derive a viable relationship between the canonical measure of \acrshort{bbo} algorithms and regret.
&  
\begin{enumerate}
    \item \textit{Aim 3.1. } Seek out the standard measures used in \acrshort{bbo} theory. 
    \item \textit{Aim 3.2. }
    \item \textit{Aim 3.2. }
    \item \textit{Aim 3.2. }
    \item \textit{Aim 3.2. }
\end{enumerate}
\\
\hline
\textbf{Aim 4. } Derive a viable relationship between the canonical measure of \acrshort{bbo} algorithms and regret.
&  
\begin{enumerate}
    \item \textit{Aim 4.1. } Seek out the standard measures used in \acrshort{bbo} theory. 
    \item \textit{Aim 4.2. }
    \item \textit{Aim 4.2. }
    \item \textit{Aim 4.2. }
    \item \textit{Aim 4.2. }
\end{enumerate}
\\
\hline
\textbf{Aim 5. } Derive a viable relationship between the canonical measure of \acrshort{bbo} algorithms and regret.
&  
\begin{enumerate}
    \item \textit{Aim 5.1. } Seek out the standard measures used in \acrshort{bbo} theory. 
    \item \textit{Aim 5.2. }
    \item \textit{Aim 5.2. }
    \item \textit{Aim 5.2. }
    \item \textit{Aim 5.2. }
\end{enumerate}
\hline
\end{longtable}


\subsection{Success Criteria}

\newpage
\section{Project Outline}
% TODO: Maybe include Gantt Chart...?
% Sem 1
%
\subsection{Literature Study}

\subsection{Isolating Relevant Benchmarks}

\subsection{Development of Experiment framework}

\subsection{Orchestrating Experiments}

\subsection{Evaluating Results}

\newpage
\section{Report Structure}
\subsection{Background}% \& Literature Review}
\subsection{Software Design}
\subsection{Experiment Structure}
\subsection{Evaluation \& Discussion of Results}
\subsection{Conclusion \& Future Work}
